\documentclass[11pt,a4paper]{article}
\usepackage[utf8]{inputenc}
\usepackage[portuguese-brazil]{babel}
\usepackage{amsmath, amssymb}
\usepackage{booktabs}
\usepackage{geometry}
\geometry{margin=2.5cm}
\usepackage{xurl}
\usepackage[colorlinks=true, linkcolor=black, urlcolor=blue, citecolor=black]{hyperref}

\title{{{title}}}
\author{Primeiro Autor \and Segundo Autor}
\date{\today}

\begin{document}
\maketitle

\begin{abstract}
Apresente o problema, a abordagem e o principal resultado em poucas frases.
\end{abstract}

\section{Introdução}

Apresente o contexto e a pergunta que este artigo responde. Cite trabalhos anteriores quando necessário \cite{knuth1984}.

\section{Método}

Descreva o que foi feito com detalhe suficiente para que outra pessoa possa repetir.

\begin{equation}
    f(x) = \frac{1}{\sigma\sqrt{2\pi}} \, e^{-\frac{(x-\mu)^2}{2\sigma^2}}
    \label{eq:normal}
\end{equation}

A Equação~\eqref{eq:normal} é referenciada automaticamente.

\section{Resultados}

\begin{table}[h]
    \centering
    \begin{tabular}{@{}lcc@{}}
        \toprule
        \textbf{Condição} & \textbf{Média} & \textbf{DP} \\
        \midrule
        Controle   & 0,00 & 1,00 \\
        Tratamento & 0,42 & 0,95 \\
        \bottomrule
    \end{tabular}
    \caption{Resumo dos resultados.}
    \label{tab:results}
\end{table}

\section{Conclusão}

Resuma os achados e suas implicações.

\begin{thebibliography}{9}
\bibitem{knuth1984}
    D.~E. Knuth, \emph{The \TeX book}. Addison-Wesley, 1984.
\end{thebibliography}

\end{document}
